\section{债务人财产}
\begin{enumerate}
  \item 破产大量考点集中在这里
  \item 非诉破产业务。破产管理人资格。
  \item 范围。到破产完成前属于它的都是。
  \item 别除权。个别的受偿。有担保的，仅有权就担保物价值获得优先受偿。登记对抗。
  \item 优先受偿。排序问题。
  \item \zd{撤销权}。参考民法债的保全。\zd{是否伤害公平，是否不合法}。
        \begin{enumerate}
          \item 转移财产、偏颇性受偿、\zd{欺诈}。管理人。\zd{1 年}内。
          \item 判断标准。是否在受理日后财产才流出、下降。
          \item \zd{针对个别清偿的撤销。破产受理前半年}（更高的注意义务）。不公平而非不合法。
          \item 个别清偿和提前清偿的叠加。
        \end{enumerate}
  \item 破产行为无效制度。逃债。
  \item 追回权。
        \begin{enumerate}
          \item 未到位的出资。
          \item 董监高非正常收入。绩效，工资高出部分、其他 \yjt 作为普通债权；工资 \yjt 职工工资清偿。
        \end{enumerate}
  \item 破产费用。立刻付，马上付。
        \begin{enumerate}
          \item 诉讼。
          \item 管理、变价、分配债务人财产产生的费用。第三方，会计、律师、拍卖行。
          \item 管理人产生的费用，应得报酬（破产分配金额比例）。
          \item 聘用工作人员费用。
        \end{enumerate}
  \item 共益债务。修复信用。立刻付，马上付。
        \begin{enumerate}
          \item 请求履行双方未履行完毕的合同产生的。
          \item 受无因管理、不当得利之债。
          \item 为继续营业而支付的报酬和社会保险费用、发生的借款。
          \item 管理人、相关人员执行职务产生的侵权。
        \end{enumerate}
  \item 清偿规则。老债老办法（普通债权），新债新办法（立即还，马上还）。
        \begin{enumerate}
          \item 随时产生，随时清偿。
          \item 顺序。破产费用 > 共益债务。但都是新债。内部按比例。
          \item 都还不起直接终结破产程序。
        \end{enumerate}
  \item 取回权。流出。权利人从债务人把东西拿回去。
        \begin{enumerate}
          \item 本质是物权的返还请求权。不是债权。
          \item 破产可以直接拿走，重整要求符合合同解除条件。
          \item 变价方案，提交债权人会议，表决之前行使。
          \item 支付运输、保管等费用 \yjt 否则我方留置权。
          \item 鲜活易腐，价款提存。
          \item 善意第三人。
          \item 原物毁损灭失。
          \item \zd{结合所有权保留买卖}。
                \begin{enumerate}
                  \item 主观题。谁破产谁说了算。继续履行 \hz 解除。
                  \item 继续履行要一次性付清全款。
                \end{enumerate}
        \end{enumerate}
  \item 抵销权。\zd{是否会造成不公平}。
        \begin{enumerate}
          \item 只能债权人向管理人主张。单向性。
          \item 没抵销掉的转为普通债权。
          \item \zd{恶意抵销}。
                \begin{enumerate}
                  \item 次债务人（债务人的债务人）破产受理后取得他人对债务人债权的。
                        \al{甲 100W \yjt 乙（破产）\yjt 丙 100W，甲出 50W 从丙手中买 100W 债权}。
                  \item 次债务人（\zd{为了抵销}）突击取得债权。
                        \al{甲欠乙 100W，卖原材料给乙，约定先欠着}。\al{超市拿糖找零钱}。
                  \item 债权人突击负债。例外：因(法律规定 \hz 破产申请 1 年前)发生的负担的债务。不当得利、无因管理、侵权。
                  \item 债权人明知债务人有不能清偿的债务 \hz 破产，对债务人负担债务。
                        \al{甲从乙拿 100W 货，约定 6 个月后付款}。
                  \item 例外。优势债权抵销劣势债权。比如本来就有别除权 \yjt 例外，不足额担保。
                \end{enumerate}
          \item 禁止股东抵销。顺序、种类、性质不一样。
        \end{enumerate}
  \item 债权申报。\zd{合法有效}。
        \begin{enumerate}
          \item 向管理人；[30 天，3 个月]；补充申报，最后分配前，费用自己承担。
          \item 未申报。除非有连带债务人 \hz 保证人，否则成为永久不能履行的权利。
          \item 老债。
          \item 罚款、不合法（高利贷）的不行。
          \item 银行 \yjt 它的债权申报 \hz 保证人 \yjt 将来求偿权申报。禁止重复申报。
          \item 保证人破产。不管主债务是否到期。丧失先诉抗辩权。分配额提存。
          \item 利息请求。
        \end{enumerate}
\end{enumerate}